\section{Uniform transfer of bounded spin functions}
\label{sec:transfer}

The scale construction requires a statement about every bounded local
function, including functions that change with the lattice scale.  We prove
that, after integration across a large square annulus, only one even scalar
direction can grow under the normalization appropriate to an interaction
density.  Convergence of any fixed collection of spin correlations would not
by itself give this assertion.  The first part of the proof supplies the
uniform approximation that permits us to use those correlations.

For an integer $s\geq1$, set
\[
 B_s=[-s,s]^2\cap\Z^2,\qquad
 B_s^\circ=\{x\in\Z^2:\|x\|_\infty<s\},\qquad
 \partial B_s=B_s\setminus B_s^\circ.
\]
Let $\mathcal X_s$ be the real functions of $\sigma_{B_s}$ invariant under
spin reversal and the dihedral group of the square, modulo constant
functions.  Its norm is
\begin{equation}
 \qnorm f:=\inf_{c\in\R}\|f-c\|_\infty
       =\tfrac12(\max f-\min f).
 \label{eq:transfer-quotient}
\end{equation}
A representative of norm $\qnorm f$ is obtained by subtracting the midpoint
of its range.  Subtracting $\E_0f$ instead gives a linear choice of
representative with norm at most $2\qnorm f$.
For an integer $L>1$, define
\begin{equation}
 A_s f=L^2\E_0[f\mid\sigma_{\partial B_{Ls}}].
 \label{eq:transfer-operator}
\end{equation}
Here the conditional expectation integrates $B_{Ls}^\circ$ at the pure
critical coupling.  Every boundary assignment has positive probability,
so this definition specifies the function on every assignment, not merely
almost surely.  It also defines a map $\mathcal X_s\to\mathcal X_{Ls}$.

\begin{theorem}[Uniform transfer]\label{thm:transfer}
There is a constant $C<\infty$ with the following property.  For each
$a>0$, all sufficiently large fixed integers $L$ admit an integer $s_0$,
linear functionals $\ell_s:\mathcal X_s\to\R$, and vectors
$e_s\in\mathcal X_s$, for every integer $s\geq s_0$, such that
\begin{align}
 \|\ell_s\|&\leq C,& \qnorm{e_s}&\leq C,& \ell_s(e_s)&=1,
 \label{eq:transfer-uniform-norms}\\
 \qnorm{A_sf}&\leq CL|\ell_s(f)|+a\qnorm f,
 \label{eq:transfer-stable}\\
 |\ell_{Ls}(A_sf)-L\ell_s(f)|&\leq a\qnorm f.
 \label{eq:transfer-compatible}
\end{align}
The constant $C$ is independent of $a,L,s$.  The choices are made in the
following order: first $L$, then a larger fixed measurement ratio $H$, then
the approximation accuracies used below, and finally $s_0$.  In particular,
all estimates hold simultaneously for all $s\geq s_0$ and all bounded
functions in $\mathcal X_s$.
\end{theorem}

We first work in fixed rescaled square charts.  A spin sum in two
dimensions is normalized by $s^{-15/8}$; the line sums appearing in
Proposition~\ref{prop:signed-transmission} are normalized by $s^{-7/8}$.
These two normalizations have different roles in the proof.
Every spatial smearing used below is a fixed bounded, compactly supported
Riemann-integrable function, chosen before the mesh limit, as in
Proposition~\ref{prop:pure-spin}.  Polygonal cell indicators belong to
this class; their specified integer roundings will be controlled below.
This convention concerns spatial averaging functions.  The microscopic
spin functions $f$ and the exterior likelihoods remain arbitrary bounded
functions.

\subsection{A window common to all exterior likelihoods}

We shall use boxes and square bands with a positive proportional buffer.
An enlargement of a band moves both boundary components outward from the
band, so its spin sets form an increasing family.  The number of boundary
components is fixed.  All distances in this subsection, unless accompanied
by a lattice scale, are in the rescaled chart.

Take such a chart with minus spins on its boundary and denote its law by
$\nu$.  Let $D$ be a layer surrounding a smaller buffered region $F_*$;
for a band it is the union of the two surrounding layers.  All altered
boundary data are farther away than $D$.  By
Proposition~\ref{prop:pure-buffered}, their marginal likelihoods on $D$
relative to $\nu$ are functions $h$ with
\begin{equation}
 0\leq h\leq K,\qquad \E_\nu h=1,
 \label{eq:transfer-layer-density}
\end{equation}
where $K$ depends only on the buffers.  Conditional expectation from $D$
to any intervening set $F$ gives its likelihood there:
$a_F=P_Fh$, where $P_F=\E_\nu[\,\cdot\mid\sigma_F]$.
These likelihoods are increasing functions of the spins of $F$.  Indeed,
relative to all-minus boundary, changing boundary spins multiplies the
joint Ising density by an increasing exponential of boundary-adjacent
spins; marginalization preserves increasing likelihood ratios for an
attractive Ising law.  This last assertion also follows directly by
conditioning on two ordered configurations of $F$ and using the monotone
coupling in the remaining spins.  Convex mixtures preserve the assertion.

\begin{lemma}[Simultaneous likelihood window]
\label{lem:transfer-window}
Fix nested enlargements $F_0\subset F_*$ of the support of a test function,
with all three relevant gaps positive: from the support to $\partial F_0$,
from $\partial F_0$ to $\partial F_*$, and from $F_*$ to $D$.
For each $t>0$ there is a finite menu of radial windows, each of positive
rescaled width, such that one window $[F_-,C]$ satisfies
\begin{equation}
 \sup_{F_-\subset F\subset C}\sup_h
       \|P_Fh-P_Ch\|_{L^2(\nu)}\leq t.
 \label{eq:transfer-window}
\end{equation}
The second supremum is over all the likelihoods in
\eqref{eq:transfer-layer-density}.  The menu and its widths are independent
of the mesh and of the boundary assignments.
\end{lemma}

\begin{proof}
Let $Q:L^2(\nu_D)\to L^2(\nu_{F_*})$ be conditional expectation from
$D$ to $F_*$.  Its integral kernel is the density of the joint law of
$(\sigma_{F_*},\sigma_D)$ relative to the product of its marginals.
The buffered comparison bounds this density by a fixed constant, hence
$\|Q\|_{\mathrm{HS}}\leq K_1$.  Choose $m+1$ equally spaced nested
enlargements $F_0\subset\cdots\subset F_m=F_*$.  Regarding all
conditional expectations as orthogonal projections in $L^2(\nu)$ gives
\[
 \sum_{j=1}^m\|(P_{F_j}-P_{F_{j-1}})Q\|_{\mathrm{HS}}^2
 \leq\|Q\|_{\mathrm{HS}}^2\leq K_1^2.
\]
Choose $m$ so that $KK_1m^{-1/2}\leq t$.  At least one summand is at
most $K_1^2/m$.  For its endpoints $F_-,C$, nesting implies, for every
intermediate $F$,
\[
 \|(P_C-P_F)h\|_2
 \leq\|(P_C-P_{F_-})Q\|_{\mathrm{op}}\|h\|_2\leq t.
\]
We used $P_Fh=P_FQh$, which holds since $F\subset F_*$.  This proves
the simultaneous, rather than likelihood-dependent, assertion.
\end{proof}

\subsection{Gaussian observations and order-positive residuals}

We now seek one observation that retains the influence of every exterior
likelihood.  Gaussian noise is useful because conditioning on it only adds
external fields to the Ising law, preserving conditional attraction.
The unobserved part of an increasing likelihood is therefore an
order-positive signed source.  If that residual stayed large, transmission
would force a magnetic signal on every one of many separating cuts.
Prediction from the observation and from nearby cuts will turn those
signals into more orthogonal increments than its bounded norm permits.

Conditional on the spins, in the selected core $C$ observe independent
Gaussian variables
\begin{equation}
 Y_x=M s^{-7/8}\sigma_x+Z_x,\qquad x\in C,
 \label{eq:transfer-observations}
\end{equation}
where the $Z_x$ are independent standard normal variables, independent
of the spins.
Write $K_Mv=\E_\nu[v\mid Y]$ and
$T_Mv=\E[K_Mv(Y)\mid\sigma]$.  Thus $T_M=K_M^*K_M$ is a
self-adjoint positive contraction on $L^2(\nu_C)$, preserves constants,
and is a contraction on bounded functions.

\begin{lemma}[Uniform smoothing of likelihoods]
\label{lem:transfer-smoothing}
Given $b>0$, the window tolerance in Lemma~\ref{lem:transfer-window}
can be chosen so that, for a sufficiently large fixed $M$ and then all
sufficiently large $s$,
\begin{equation}
 \sup_h\|(I-T_M)P_Ch\|_2<b.
 \label{eq:transfer-smoothing}
\end{equation}
All constants are uniform over the likelihood family and the finite menu
of possible windows.
\end{lemma}

\begin{proof}
Put $a=P_Ch$ and $z=(I-T_M)a$.  Then $\E_\nu z=0$, $|z|\leq2K$,
and, for every increasing $g$ of the core spins,
\begin{equation}
 \E_\nu[zg]=\E[\Cov_\nu(a,g\mid Y)]\geq0.
 \label{eq:transfer-order-positive}
\end{equation}
To justify the last inequality, conditioning on $Y=y$ changes the Ising
weight by $\exp(Ms^{-7/8}\sum_x y_x\sigma_x)$ times a constant.
The conditional law remains ferromagnetic, with external fields of
arbitrary signs, so conditional FKG applies
\citep[Proposition~1 and pp.~98--99]{FKG1971}.  The function $z$ itself
need not be increasing.  Positivity of $I-T_M$ and its contraction bound
also give
\begin{equation}
 \langle z,a\rangle
 =\langle a,(I-T_M)a\rangle\geq\|z\|_2^2.
 \label{eq:transfer-residual-energy}
\end{equation}

Suppose $\|z\|_2\geq b$.  Choose the window tolerance less than
$b^2/(4K)$.  For every $F$ in that window,
\[
 \langle P_Fz,h\rangle=\langle z,P_Fh\rangle
 \geq b^2-2Kt\geq b^2/2.
\]
Subdivide an inner fixed portion of the window into $n$ equal radial
steps.  Denote the successive spin sets by $F_0',\ldots,F_n'$ and choose
a separating square contour, or the two contours of a band, $J_i$
strictly between $F_i'$ and $F_{i+1}'$.  Its distance from $F_i'$ is
comparable to $w=s/n$, with a constant depending on the already fixed
window.  Define $m_i=s^{-7/8}\sum_{x\in J_i}\sigma_x$.
The signed source $P_{F_i'}z$ is order-positive: an increasing
$F_i'$-measurable test has the same pairing with it as with $z$.
Subdivide $J_i$ into straight pieces of a sufficiently small fixed
multiple of $w$.  Their enlarged charts are pin-free, they have bounded
overlap, and the target layer $D$ is at distance comparable to $s$.
These are exactly the admissible cuts of
Proposition~\ref{prop:signed-transmission}.  For its fixed exponent
$\eta>0$, that proposition yields
\begin{equation}
 \langle P_{F_i'}z,m_i\rangle\geq c_b n^{-7/8+\eta}.
 \label{eq:transfer-transmission-lower}
\end{equation}
The constant does not deteriorate further as $n$ increases.

The observation predicts each $m_i$ with arbitrarily small error when
$n$ is fixed.  To see this quantitatively, average it over parallel
nearby lines, including the corresponding short endpoint corrections,
in a strip of width $\theta s$ lying inside $C$.  Call the average
$\bar m_i$.  All these lines have the same length and buffer bounds,
with displacement at most $C\theta s$; a square contour is treated
side by side, including its short endpoint pieces.  For fixed $\theta$,
Proposition~\ref{prop:line-continuity} is uniform over the averaging
lines: otherwise choose a maximizing line at each mesh, obtaining a
sequence of pairs that still satisfies that proposition with
$w/s=\theta$.  Minkowski's inequality therefore gives, even though
the number of averaging lines grows with $s$,
\[
 \limsup_{s\to\infty}\|m_i-\bar m_i\|_2\leq C\theta^{3/8}.
\]
The average is a linear combination of the signals
$s^{-7/8}\sigma_x$, with coefficients of size $O((\theta s)^{-1})$
on $O(\theta s^2)$ sites.  Replacing each signal by $Y_x/M$ gives a
$Y$-measurable estimator of $\bar m_i$ with noise variance at most
$C/(M^2\theta)$.  Consequently
\[
 \limsup_{s\to\infty}\|m_i-\E[m_i\mid Y]\|_2
 \leq C\theta^{3/8}+C/(M\sqrt\theta).
\]
Since
$\langle z,m_i\rangle=\E[\Cov(a,m_i\mid Y)]$, first choosing
$\theta$ small and then $M$ large makes this pairing smaller in absolute
value than $\tfrac12c_b n^{-7/8+\eta}$.

Set $d_i=(P_{F_{i+1}'}-P_{F_i'})z$.  The variable $m_i$ is measurable
on $F_{i+1}'$, whereas $d_i$ is orthogonal to all $F_i'$-measurable
functions.  Therefore
\[
 |\langle d_i,m_i-P_{F_i'}m_i\rangle|
 =|\langle z,m_i\rangle-\langle P_{F_i'}z,m_i\rangle|
 \geq\tfrac12c_b n^{-7/8+\eta}.
\]
Corollary~\ref{cor:line-prediction}, with a nearby contour in $F_i'$
and with $n$ fixed before $s\to\infty$, bounds its second factor by
$Cn^{-3/8}$.  Hence $\|d_i\|_2\geq c_b'n^{-1/2+\eta}$ for all
sufficiently large $s$.  Orthogonality of the increments gives the
contradiction
\begin{equation}
 4K^2\geq\|z\|_2^2\geq\sum_{i=0}^{n-1}\|d_i\|_2^2
             \geq c_b''n^{2\eta}.
 \label{eq:transfer-contradiction}
\end{equation}
Choose $n$ large enough before choosing $\theta,M,s$.  Thus the complete
order is $b$, the window tolerance and width, $n$, $\theta$, $M$, and
finally the mesh threshold.  Only finitely many contours and windows
occur after $n$ has been fixed, so their thresholds can be combined.
\end{proof}

This gives an approximation in exactly the norm needed for separated
correlations.  If $v$ is a bounded function on the original support,
then for every exterior likelihood,
\begin{equation}
 |\E_\nu[(v-T_Mv)h]|
 =|\langle v,(I-T_M)P_Ch\rangle|
 \leq\|v\|_\infty b.
 \label{eq:transfer-testing-error}
\end{equation}
The same conclusion holds against a bounded exterior observable, by
conditioning on its spin configuration and integrating the preceding
bound.  The remaining issue is that $Y$ has a number of coordinates
growing with $s$.  We next reduce it to a fixed finite collection.

\subsection{Overlap moments and finitely many observation modes}

For two independent copies $\sigma,\sigma'$ on a bounded rescaled
region, put
\[
 Q_s=s^{-7/4}\sum_x\sigma_x\sigma_x'.
\]
We shall use the following two estimates, also valid under any pure law
whose density in the region is bounded relative to the plane law:
\begin{align}
 \sup_s\E\exp(t|Q_s|)&<\infty\quad(t<\infty),
 \label{eq:transfer-overlap-exponential}\\
 \lim_{\kappa\downarrow0}\limsup_{s\to\infty}
   \E|Q_s-Q_{s,\kappa}|^2&=0.
 \label{eq:transfer-overlap-projection}
\end{align}
Here $\Pi_\kappa$ is orthogonal projection of site vectors onto functions
constant on the cells of a fixed rectangular partition of mesh $\kappa$
in rescaled coordinates, intersected with the selected square or
square-band core.  These cells are polygonal and have boundaries of area
zero.  Retain the positive-area cells, assigning boundary sites to one
adjacent cell by a fixed rule.  The finite partition is fixed before
the mesh limit.  Define
$Q_{s,\kappa}=s^{-7/4}\langle\Pi_\kappa\sigma,
\Pi_\kappa\sigma'\rangle$.

We now control the growth in the moment order to prove the uniform
integrability needed in these statements.  Recall~\eqref{eq:pure-multipoint-bound}:
for $n\geq2$ and sites that need not be distinct,
\begin{equation}
 \left|\E_0\prod_{i=1}^n\sigma_{x_i}\right|
 \leq C^n\prod_{i=1}^n
       (1+\min_{j\ne i}|x_i-x_j|)^{-1/8}.
 \label{eq:transfer-nearest-moment}
\end{equation}
Expanding an even moment of $Q_s$ squares the correlation in
\eqref{eq:transfer-nearest-moment}.  After rescaling, it remains to
bound the integral, with normalized lattice counting measure on a fixed
bounded region, of
$\prod_i(\delta\vee r_i)^{-1/4}$, where $\delta=s^{-1}$ and
$r_i=\min_{j\ne i}|x_i/s-x_j/s|$.  Take $h=n^{-1/2}$.
If $\delta>h$, the deterministic bound $|Q_s|\leq Cs^{1/4}$ suffices.
Otherwise let
$\rho_i=\min\{h,\max\{\delta,r_i\}\}$, and randomly order the
$n$ labels.  Write $D_i$ for the distance to the nearest predecessor,
with the same cutoffs and with $D_i=h$ if there is no predecessor.
A selected nearest neighbor precedes label $i$ with probability $1/2$.
Thus Jensen's inequality gives
\begin{equation}
 \prod_i\rho_i^{-1/4}
 \leq h^{-n/4}\E_{\mathrm{order}}
        \prod_i(h/D_i)^{1/2}.
 \label{eq:transfer-ordering}
\end{equation}
For any fixed ordering, integrate positions in reverse order.  For each
last remaining position the union of disks of radius $r\geq\delta$
around its predecessors has normalized counting measure at most
$Cnr^2$.  The layer-cake integral for $(h/D_i)^{1/2}$ is bounded by a
constant, since $nh^2=1$; the atom at the cutoff $\delta$ obeys the same
bound.  Repeating the integration proves
\begin{equation}
 \E_0 Q_s^n\leq C^n n^{n/8}\quad(n\text{ even}).
 \label{eq:transfer-overlap-moments}
\end{equation}
Cauchy--Schwarz gives the corresponding absolute-moment bound for all
integer orders.  Its exponential series converges for every fixed
coefficient, proving \eqref{eq:transfer-overlap-exponential}.

For \eqref{eq:transfer-overlap-projection}, the identity for the
second moment of an overlap expresses the error through the covariance
operator.  More explicitly, with $C_s$ the covariance matrix of
$s^{-7/8}(\sigma_x)_x$ in ordinary Euclidean coordinates,
\[
 \E_0|Q_s-Q_{s,\kappa}|^2
 =\operatorname{tr}\left(((I-\Pi_\kappa)C_s
                     (I-\Pi_\kappa))^2\right)
 \leq\|C_s-\Pi_\kappa C_s\Pi_\kappa\|_{\mathrm{HS}}^2.
\]
The two-point convergence and its bound
$C(1+|x-y|)^{-1/4}$ identify the limiting kernel with a constant times
$|u-v|^{-1/4}$.  Its square is locally integrable in two dimensions.
Removing a small neighborhood of the diagonal, taking Riemann sums,
and then restoring that neighborhood proves Hilbert--Schmidt convergence
and the claim.  Buffered density bounds transfer the nonnegative error
and the exponential estimates to the reference laws.

The coarse overlaps also have the exponential bounds uniformly in
$\kappa$.  Under the plane law, if
$m_n=\E[(s^{-7/8}\sigma)^{\otimes n}]$, then
\[
 \E Q_{s,\kappa}^n
 =\|\Pi_\kappa^{\otimes n}m_n\|^2\leq\|m_n\|^2
 =\E Q_s^n\qquad(n\text{ even}).
\]
Again density comparison gives the assertion in a buffered reference
chart.  It follows from \eqref{eq:transfer-overlap-projection}, truncation,
and these exponential bounds that, for fixed $M$,
\begin{equation}
 \E\left|e^{M^2Q_s}-e^{M^2Q_{s,\kappa}}\right|\longrightarrow0
 \quad\text{as }s\to\infty\text{ and then }\kappa\downarrow0.
 \label{eq:transfer-exponential-projection}
\end{equation}

We now apply this calculation to observations.  Relative to standard
white noise $\gamma_s$ on the sites of $C$, the observation density
weighted by a function $v$, $|v|\leq1$, is
\[
 R_v(y)=\E_\nu\left[v(\sigma)
  \exp\{M\langle y,s^{-7/8}\sigma\rangle
        -\tfrac12M^2\|s^{-7/8}\sigma\|^2\}\right].
\]
Let $R_v^\kappa=\E_{\gamma_s}[R_v\mid\Pi_\kappa y]$.  Gaussian
integration and the orthogonal-projection identity give exactly
\begin{equation}
 \|R_v-R_v^\kappa\|_{L^2(\gamma_s)}^2
 =\E_{\nu\otimes\nu}\left[v(\sigma)v(\sigma')
      (e^{M^2Q_s}-e^{M^2Q_{s,\kappa}})\right].
 \label{eq:transfer-output-density}
\end{equation}
The right side tends to zero uniformly over $|v|\leq1$.
The full and coarse posterior means are $R_v/R_1$ and
$R_v^\kappa/R_1^\kappa$.  They are bounded by one, and
\begin{equation}
 \int R_1\left|\frac{R_v}{R_1}
       -\frac{R_v^\kappa}{R_1^\kappa}\right|d\gamma_s
 \leq\|R_v-R_v^\kappa\|_1+\|R_1-R_1^\kappa\|_1.
 \label{eq:transfer-ratio}
\end{equation}
This uses $|R_v^\kappa|\leq R_1^\kappa$ and requires no positive
uniform lower bound on $R_1$.

Write $X_{s,\kappa}=s^{-7/8}\Pi_\kappa\sigma$ in an orthonormal
cell basis, and let $q_{s,v}$ be the coarse posterior as a function of
the projected observation.  Its absolute value is at most one.
Conditioning back on the spins gives the replacement
\[
 F_{s,v}(X_{s,\kappa}),\qquad
 F_{s,v}(x)=\int q_{s,v}(Mx+z)\,d\gamma_{d_\kappa}(z),
\]
where $d_\kappa$ is the fixed number of retained cells and
$\gamma_{d_\kappa}$ is standard Gaussian measure in that dimension.
Thus it is the Gaussian convolution of the bounded posterior that is
uniformly Lipschitz, with a constant depending on $M$ and the partition
but not on $v$ or $s$: bound the difference by the total variation
distance between the two translated Gaussians.
For a fixed positive-area polygonal cell
$E$, put $n_{s,E}=|sE\cap C|$, with the prescribed boundary assignment.
In its normalized cell basis the corresponding coordinate is
\[
 s^{-7/8}n_{s,E}^{-1/2}\sum_{x\in sE\cap C}\sigma_x
 =\left(\frac{s^2}{n_{s,E}}\right)^{1/2}
       s^{-15/8}\sum_{x\in sE\cap C}\sigma_x.
\]
Here $n_{s,E}/s^2\to\operatorname{area}(E)>0$, so the cell is nonempty
for all sufficiently large $s$ and the scalar normalization converges.
Its spatial smearing is the fixed Riemann-integrable indicator of $E$.
Equation~\eqref{eq:transfer-ratio} gives an $L^1(\nu)$
error, which is sufficient against every uniformly bounded likelihood.

\begin{lemma}[Bounded replacement across a gap]
\label{lem:bounded-replacement}
Fix finitely many pairs of positively separated supports, each a square
or a square band, with buffered enlargements that remain disjoint.
For every $t>0$, each lattice function $v$ of norm at most one on any
one support can be replaced, for its full-plane correlations against all
functions of norm at most one on the paired supports, with error at most
$t$, by a function
of finitely many smeared spin sums in its enlargement.  The replacement
has norm at most one.  The number of sums, their shapes, and a modulus
of continuity for the replacing functions belong to a fixed finite
menu, independent of $v$ and of all sufficiently fine meshes.
Along any mesh sequence, further subsequences of these replacements
converge to bounded continuous tests of finitely many continuum smeared
spins.  Replacements on both sides of each gap can be made together,
and preserve spin-flip and square symmetries by averaging.
\end{lemma}

\begin{proof}
Use Lemma~\ref{lem:transfer-smoothing}, then
\eqref{eq:transfer-testing-error}, and finally the finite-mode
construction.  To pass from the reference law $\nu$ to a full-plane
pairing, choose the buffered reference chart inside the specified
enlargement, disjoint from the paired supports, and condition the plane
law on its exterior spins.  The
conditional marginal on $D$ has a likelihood $h$ in
\eqref{eq:transfer-layer-density}, and its core likelihood is $P_Ch$.
The paired observable is measurable outside that chart, so integrating the bound
against its signed values gives the same correlation estimate under
$\mu_0$.  For the finite-mode error, every such core likelihood is
bounded by $K$; its $L^1(\nu)$ error therefore costs at most that fixed
factor.  The replacing function is the same for every exterior
configuration.  For finitely many required supports the menu is
their finite product, and errors are divided by their number.  The
bounded functions of the finite-dimensional sums are equicontinuous;
tightness of those sums and compactness on bounded coordinate sets give
the subsequential assertion by a diagonal argument.  These replacements
and their smeared coordinates are now evaluated under $\mu_0$.
Their limiting field law is therefore the full-plane spin-field law:
Proposition~\ref{prop:pure-spin} gives joint moment convergence for the
fixed polygonal cell indicators.  The
prescribed rounding errors lie in vanishing neighborhoods of their
finitely many edges.  The uniform $L^2$ bounds on fixed-order rescaled
correlation kernels proved there, and the vanishing lattice counting
measure of these neighborhoods, make the errors tend to zero by
Cauchy--Schwarz.  The finite menu of spatial smearings is fixed before
the mesh limit.  The unsquared version of the ordering argument
\eqref{eq:transfer-ordering} supplies exponential integrability of each
fixed smeared sum and hence moment determinacy.  Averaging over the
finite symmetry group preserves all bounds.  Applying the result to
one side and then the other costs at most the sum of the two errors.
\end{proof}

The lemma asserts approximation for correlations across a gap.  No
pointwise recovery of an arbitrary spin function from the magnetization
field is asserted or needed.

\subsection{The radial Hilbert space}

We have reduced lattice correlations of bounded functions to continuum
field tests with controlled ordinary $L^2$ norm.  We now identify the
part of their radial spectrum below degree three.  The explicit spin
formula will first give expansions for polynomial tests; reflection
positivity will then turn those expansions into an operator bound.

Use the normalization of Proposition~\ref{prop:pure-spin}, so the
continuum two-point function is $|z-w|^{-1/4}$, and write $\Phi(\psi)$
for the spin field smeared against a bounded compactly supported
Riemann-integrable spatial function $\psi$.  All spatial smearings in
this section belong to this class.  Let $\mathcal A$ be the algebra of
polynomials in such smeared fields supported strictly inside the unit
disk.  Initially their supports may be kept away from the origin.
This causes no loss upon completion: multiply each spatial function by
a smooth cutoff vanishing on a disk of radius $\varepsilon$ and equal to one
outside the disk of radius $2\varepsilon$.  As $\varepsilon\downarrow0$
these functions converge almost everywhere, remain uniformly bounded,
and have common compact support.  The integrable fixed-order correlation
bounds give convergence of each smeared variable in every fixed finite
$L^p$.  Telescoping a polynomial difference and applying H\"older's
inequality therefore gives ordinary $L^2$ convergence of every fixed
smeared polynomial.

The spatial class is preserved by the weighted coordinate changes used
below, when the support stays away from their poles.  If $g$ is a
M\"obius or anti-M\"obius map, denote its local length scale by
$|Dg(z)|$, equal to $|g'(z)|$ in the conformal case, and set
$T_g\sigma(z)=|Dg(z)|^{1/8}\sigma(g(z))$ on point insertions.
Changing variables in a smearing gives
\begin{equation}
 T_g\Phi(\psi)=\Phi(\psi_g),\qquad
 \psi_g(w)=\psi(g^{-1}(w))|Dg^{-1}(w)|^{15/8}.
 \label{eq:transfer-spatial-transform}
\end{equation}
The area Jacobian has exponent $2$ and the primary spin weight subtracts
$1/8$, giving the exponent $15/8$ in this formula.  On the relevant
compact neighborhoods the coordinate map and its inverse are smooth
with bounded derivatives, preserve null sets, and have a bounded
continuous Jacobian weight.  Thus $\psi_g$ is bounded, compactly
supported, and Riemann integrable.  In particular inversion gives the
spatial weight $|w|^{-15/4}\psi(1/\bar w)$.

Cutoffs are made before applying a map with a pole.  For a polynomial
$F$ in the cutoff smears, correlation covariance gives
$\|T_gF\|_2=\|F\|_2$.  After removing the cutoffs, this isometry
defines the transformed polynomial by an ordinary $L^2$ limit.  For
example, inverted supports need not remain in one compact set as an
origin cutoff is removed.  The extension uses the isometry and asserts
no lattice sampling theorem for the resulting unbounded-support
expression.

Polynomials are dense in the $L^2$ space of any finite collection of
smeared fields.  Here is the determinacy argument in the precise form
needed.  Apply the random predecessor ordering argument of
\eqref{eq:transfer-ordering} to the unsquared nearest-distance bound,
with exponent $1/8$ in place of $1/4$.  For even $n$ it gives
\[
 \E|\Phi(\psi)|^n\leq C_\psi^n n^{n/16}.
\]
Cauchy--Schwarz gives the same bound, after enlarging $C_\psi$, for
every integer $n\geq1$.  Thus $\E e^{t|\Phi(\psi)|}<\infty$ for
every fixed $t$.  If an $L^2$
function $g$ of a finite vector of smears is orthogonal to all its
polynomials, its signed distribution has a Laplace transform analytic in
a neighborhood of the origin, by Cauchy--Schwarz and exponential
integrability.  Every derivative there is zero.  Analytic continuation
to the imaginary axes and uniqueness of Fourier transforms give $g=0$.
Thus bounded continuous tests in Lemma~\ref{lem:bounded-replacement}
can be approximated in ordinary $L^2$ by elements of $\mathcal A$.

We use the reflection-positive Hilbert-space construction of
\citet[Section~4.1]{OsterwalderSchrader1973} in its radial form,
verifying contraction and continuity below.
Let $\iota(z)=1/\bar z$.  On point insertions, interpreted through
smearing, reflection in the unit circle acts as
\[
 \Theta\sigma(z)=|z|^{-1/4}\sigma(\iota(z)).
\]
It acts multiplicatively on polynomials.  Define
\begin{equation}
 (F,G)_{\mathrm{rad}}=\E[(\Theta\overline F)G],
 \qquad \|F\|_{\mathrm{rad}}^2=(F,F)_{\mathrm{rad}}.
 \label{eq:transfer-radial-form}
\end{equation}
The field is invariant under this transformation by its M\"obius
covariance.  Line reflection positivity from
Lemma~\ref{lem:pure-reflection}, transported to the circle by a
M\"obius map, makes this form positive semidefinite.  Cauchy--Schwarz
in the ordinary probability space gives
\begin{equation}
 \|F\|_{\mathrm{rad}}\leq\|F\|_2.
 \label{eq:transfer-radial-l2}
\end{equation}
Let $\mathcal H$ be its completion after quotienting out the null space.

For $t\geq0$, let $D_t$ dilate insertions inward by $r=e^{-t}$, including
their primary weights:
$D_t\sigma(z)=r^{1/8}\sigma(rz)$.
The covariance formula shows that $D_tD_u=D_{t+u}$, that $D_t$ is
symmetric for \eqref{eq:transfer-radial-form}, and that
$\|D_tF\|_2=\|F\|_2$.  To verify the needed contraction without
assuming a spectral representation, use symmetry and Cauchy--Schwarz:
\[
 \|D_tF\|_{\mathrm{rad}}^2
 =(F,D_{2t}F)_{\mathrm{rad}}
 \leq\|F\|_{\mathrm{rad}}\|D_{2t}F\|_{\mathrm{rad}}.
\]
Iteration, followed by \eqref{eq:transfer-radial-l2}, gives
\[
 \|D_tF\|_{\mathrm{rad}}
 \leq\|F\|_{\mathrm{rad}}^{1-2^{-k}}
       \|F\|_2^{2^{-k}}.
\]
Letting $k\to\infty$ proves contraction, and proves as well that
$D_t$ descends through the null space.  For $r\to1$, the spatial
smearing associated with $D_{-\log r}$ is
$r^{-15/8}\psi(\,\cdot/r)$.  It converges to $\psi$ at every
continuity point of $\psi$, hence almost everywhere, and the functions
are uniformly bounded on a common compact set.  Dominated convergence
against the integrable fixed-order correlation bounds gives convergence
of each smeared variable in every fixed finite $L^p$.  Telescoping
polynomial differences and using H\"older's inequality gives ordinary
$L^2$ convergence of every fixed smeared polynomial.  By
\eqref{eq:transfer-radial-l2} this proves strong continuity on
$\mathcal A$, then on $\mathcal H$ by contraction. The self-adjoint
contraction-semigroup theorem
\cite[Theorem~7.3.10, pp.~337--338]{BuehlerSalamon2016} therefore gives
\begin{equation}
 D_t=e^{-t\mathcal D},\qquad \mathcal D\geq0,
 \label{eq:transfer-radial-generator}
\end{equation}
with $\mathcal D$ self-adjoint. The cited real-Hilbert-space theorem
applies to the underlying real space; complex linearity of the semigroup
makes its generator complex linear as well. Square symmetries and spin reversal
commute with this semigroup.

\subsection{The complex determinant identity and the scalar gap}

\begin{lemma}[Charged sums at complex points]
\label{lem:transfer-charges}
Let $z_1,\ldots,z_{2m}$ be distinct complex points, with $m\geq1$, and
put
\[
 W_z(u)=\prod_{i<j}|z_i-z_j|^{u_i u_j/2},\qquad
 A_q(z)=\sum_{\substack{u_i\in\{-1,1\}\\\sum_i u_i=q}}W_z(u),
\]
where an empty charge class has sum zero.  Define
\[
 A_0^d(z)=\sum_{\sum_i u_i=0}W_z(u)
                        \left|\sum_i u_i z_i\right|^2.
\]
Then
\begin{equation}
 A_0^dA_0=4A_2^2.
 \label{eq:transfer-charge-identity}
\end{equation}
If all points lie in the unit disk, then $A_2/A_0\leq1$, and all
ratios $A_q/A_0$ are bounded by one.
\end{lemma}

\begin{proof}
Set $w_i=(\prod_{j\ne i}|z_i-z_j|)^{-1}$ and
$C_z=\prod_{i<j}|z_i-z_j|^{1/2}$.  For the positive subset
$I=\{i:u_i=1\}$, direct multiplication gives
\begin{equation}
 W_z(u)=C_z\,|\Delta(z_I)|^2\prod_{i\in I}w_i,
 \label{eq:transfer-vandermonde}
\end{equation}
where $\Delta$ is the Vandermonde determinant.  Let $Z_k$ be the sum of
the determinant weights on the right, without $C_z$, over $|I|=k$;
set $Z_0=1$.  For the positive measure $\sum_iw_i\delta_{z_i}$ on
this finite set of complex points,
let $p_j$ be the monic orthogonal polynomial of degree $j$, and let
$h_j=\sum_iw_i|p_j(z_i)|^2$.  Cauchy--Binet applied to the evaluation
matrix gives
\begin{equation}
 Z_k=h_0\cdots h_{k-1},\qquad 1\leq k\leq2m.
 \label{eq:transfer-gram}
\end{equation}
The weight in \eqref{eq:transfer-vandermonde} without $C_z$ is also
$\prod_{i\in I,j\notin I}|z_i-z_j|^{-1}$.
It is unchanged under complementation.  Thus $Z_{2m-k}=Z_k$ and
\begin{equation}
 h_mh_{m-1}=Z_{m+1}/Z_{m-1}=1.
 \label{eq:transfer-complement}
\end{equation}

Consider the probability distribution on $m$-subsets proportional to
these weights.  Its inclusion probabilities are the minors of the
orthogonal projection $P$ onto polynomials of degree at most $m-1$,
viewed on the finite weighted point set.  This is another direct
Cauchy--Binet identity: summing squared minors containing a prescribed
set leaves the corresponding principal minor of $P$.
Let $M_z$ denote multiplication by $z$ on that space.  Expanding the
absolute-square variance with the one- and two-point inclusion
probabilities gives
\begin{align*}
 \Var\left(\sum_{i\in I}z_i\right)
 &=\operatorname{tr}(PM_z^*M_z)
       -\operatorname{tr}(PM_z^*PM_z)\\
 &=\|(I-P)M_zP\|_{\mathrm{HS}}^2
   =h_m/h_{m-1}.
\end{align*}
The last equality uses the orthonormal polynomial basis: multiplication
by $z$ sends its first $m-1$ vectors into $\operatorname{ran}P$, and
the last vector has complement equal to $p_m/\sqrt{h_{m-1}}$.
This argument uses complex absolute-square variance; no assumption that
the points lie on a real line is present.

Complementation shows that
$\E\sum_{i\in I}z_i=\tfrac12\sum_i z_i$.  Since
$\sum_i u_i z_i=2\sum_{i\in I}z_i-\sum_i z_i$, we obtain
\[
 \frac{A_0^d}{A_0}=4\frac{h_m}{h_{m-1}}=4h_m^2,
 \qquad \frac{A_2}{A_0}=\frac{Z_{m+1}}{Z_m}=h_m,
\]
which proves \eqref{eq:transfer-charge-identity}.
If $|z_i|\leq1$, the minimizing property of monic orthogonal
polynomials, applied to $zp_{j-1}$, gives $h_j\leq h_{j-1}$.
Together with \eqref{eq:transfer-complement} this implies $h_m\leq1$.
Ratios $Z_{m+j}/Z_m=h_m\cdots h_{m+j-1}$ are at most one, and
complementation treats negative charges.
\end{proof}

\begin{proposition}[Radial spectral gap]\label{prop:radial-gap}
On the subspace of $\mathcal H$ generated by even, quarter-turn
invariant spin tests, the spectrum of $\mathcal D$ below $3$ consists
of the constant state at $0$ and one nonzero one-dimensional state at
$1$.  There is a real linear functional $\ell$ with
$|\ell(F)|\leq\|F\|_2$ such that, for real even quarter-turn
invariant tests $F,G$ of finite ordinary $L^2$ norm and $0<r\leq1$,
\begin{equation}
 \left|\E[(\Theta F)D_{-\log r}G]
       -\E F\E G-r\ell(F)\ell(G)\right|
 \leq r^3\|F\|_2\|G\|_2.
 \label{eq:transfer-radial-gap}
\end{equation}
The assertion extends from smeared polynomials to bounded field tests
by ordinary $L^2$ approximation.
\end{proposition}

\begin{proof}
We compute first with an even number $n=2m$ of inner insertions $z_i$
and $n'=2m'$ reflected outer insertions $y_j$.  Both insertion sets
are inside the unit disk and have distinct points.  After including
the dilation and reflection weights, the squared spin formula in
Proposition~\ref{prop:pure-spin} becomes
\begin{equation}
 2^{-(n+n')/2}\sum_q r^{q^2/4}
 \sum_{\substack{\sum u_i=q\\\sum v_j=-q}}
 W_z(u)W_y(v)
 \prod_{i,j}|1-rz_i\bar y_j|^{u_i v_j/2}.
 \label{eq:transfer-radial-expansion}
\end{equation}
For example, the power of $r$ is the primary factor $r^{n/4}$ in the
squared expression times the internal-distance factor
$r^{(q^2-n)/4}$.  All powers of $|y_j|$ from inversion cancel.
This verifies the normalization of the radial degree.

Angular frequency refers here to the terms $e^{ij\theta}$ produced by
rotating one whole insertion set, $z_i\mapsto e^{i\theta}z_i$.
Divide \eqref{eq:transfer-radial-expansion} by its constant term,
and abbreviate the charged sums of the $y$ variables by $B_q,B_0^d$.
The $q=\pm2$ terms have first-order coefficient
\[
 p=2(A_2/A_0)(B_2/B_0).
\]
In the charge-zero term the first-order cross factor is proportional
to
$\operatorname{Re}[(\sum_i u_i z_i)
\overline{(\sum_jv_j y_j)}]$ and vanishes after summing either sign
family.  To find the angularly constant second-order coefficient,
use
\[
 \log\prod_{i,j}|1-rz_i\bar y_j|^{u_iv_j/2}
 =-\tfrac r2\operatorname{Re}(X\bar Y)
   -\tfrac{r^2}{4}\operatorname{Re}(X_2\overline{Y_2})+O(r^3),
\]
where $X=\sum u_i z_i$, $X_2=\sum u_i z_i^2$, and similarly for
$Y,Y_2$.  Squaring the linear term in the exponential contributes
$r^2|X|^2|Y|^2/16$ in angular frequency zero.  Its normalized sum is
\[
 q_0=\frac{A_0^dB_0^d}{16A_0B_0}.
\]
Every other second-order term has angular frequency one or two: frequency
one comes from the first cross correction in the $q=\pm2$ sector, and
frequency two from the remaining charge-zero terms.  Charges of absolute
value at least four begin at $r^4$.

The correlation itself is the positive square root.  Its scalar
second-order coefficient is $q_0/2-p^2/8$.  By
Lemma~\ref{lem:transfer-charges}, $q_0=p^2/4$, so this coefficient
is zero.  Averaging one insertion set over quarter turns kills the
nonzero angular frequencies just listed.  The first-order coefficient
in the unsquared correlation is
\begin{equation}
 \bigl(S(z_1,\ldots,z_n)A_2/A_0\bigr)
 \bigl(S(y_1,\ldots,y_{n'})B_2/B_0\bigr).
 \label{eq:transfer-factorized-coefficient}
\end{equation}
It factors across the gap.  An empty insertion set contributes only
the constant state.

These computations survive smearing.  For a fixed number of insertions
and $|r|$ sufficiently small, expand the real-analytic factors into
their absolutely convergent power series in $r$ and the angular
variables.  The ratios of charged sums are bounded by
Lemma~\ref{lem:transfer-charges}; the cross factors and their first
three derivatives are bounded uniformly when $|z_i|,|y_j|\leq1$ and
$|r|\leq r_0<1$.  The remaining integrable dominating function is a
constant times the product of the within-set spin correlations.
Thus every even quarter-turn invariant smeared polynomial has the
bilinear expansion
\begin{equation}
 (F,D_{-\log r}G)_{\mathrm{rad}}
   =\E\overline F\E G+r\overline{c(F)}c(G)+O_{F,G}(r^3),
 \label{eq:transfer-polynomial-expansion}
\end{equation}
with the linear functional $c$ given by
\eqref{eq:transfer-factorized-coefficient} and linear extension.

We explain why the polynomial-dependent remainder in this calculation
implies the uniform bound in the proposition.  In a self-pairing the
spectral functional calculus
\cite[Theorem~6.4.1 and Section~6.5, especially Theorem~6.5.4]{BuehlerSalamon2016}
represents the left side of \eqref{eq:transfer-polynomial-expansion} as
$\int_{[0,\infty)}r^d\,d\mu_F(d)$ with $\mu_F$ positive.
Letting $r\downarrow0$ identifies its mass at zero as $|\E F|^2$.
After subtracting it, divide by $r$.  Positivity excludes mass on
$(0,1)$ and identifies the mass at $1$ as $|c(F)|^2$.
After subtracting this atom as well, the remainder is positive and
$O_F(r^3)$, which excludes mass on $(1,3)$: any positive mass in a
compact interval with upper endpoint below $3$ would violate that
bound.  Density of the polynomial tests excludes these spectral
intervals on the entire stated subspace.

Polarization and the two factorizations in
\eqref{eq:transfer-polynomial-expansion} show that the projections at
$0$ and $1$ have rank one, provided the latter is nonzero.  It is
nonzero already for a two-spin test smeared in two small disjoint
regions: choose nonzero nonnegative smooth compactly supported bumps
$\psi_1,\psi_2$ in those regions and use
$\Phi(\psi_1)\Phi(\psi_2)$.  The coefficient is strictly positive
since $S(z_1,z_2)A_2/A_0>0$ for distinct points; square-group
averaging preserves this positive scalar coefficient.  Choose a unit real vector in that
eigenspace and denote the corresponding coefficient by $\ell$.
Then $|\ell(F)|\leq\|F\|_{\mathrm{rad}}\leq\|F\|_2$.
On the remaining spectral subspace, the norm of $D_{-\log r}$ is at
most $r^3$.  Cauchy--Schwarz proves
\eqref{eq:transfer-radial-gap}.  Finally
\eqref{eq:transfer-radial-l2} and polynomial density give the asserted
extension.
\end{proof}

In particular there exist bounded continuous, even, square-invariant
tests with a nonzero degree-one coefficient.  Start with the smeared
two-spin test built from the smooth bumps in the proof and average it
under the square group.  Continuous bounded truncations converge to
this test in ordinary $L^2$ and are functions of the same finite vector
of smooth smears.  Since the coefficient is continuous in this norm,
a sufficiently large fixed truncation has nonzero coefficient.
Choose the bumps in a fixed annular subset of the disk.  Their square
images remain in that annulus, so reflection keeps all transformed
spatial smears smooth and compactly supported in a bounded annulus.

\subsection{Radial coefficients for square supports}

We first record how to use Proposition~\ref{prop:radial-gap} for square
supports.  If a continuum test $F$ is supported in a fixed disk of
radius $a_0$, define its degree-one coefficient by
\[
 \Lambda(F)=a_0\ell(D_{\log a_0}F),
\]
taking $a_0\geq1$ large enough that the dilated support is strictly
inside the unit disk.  This definition does not depend on $a_0$:
the unit degree eigenstate scales by $e^{-t}$ under $D_t$.
In particular $|\Lambda(F)|\leq C\|F\|_2$ on any fixed bounded
family of supports.  For a test $G$ supported in a fixed annulus with
inner radius $b_0>0$, choose a slightly smaller positive $b<b_0$ and
put
\[
 \Gamma(G)=b^{-1}\ell(\Theta D_{\log b}G).
\]
Here a dilation with a negative parameter is simply the primary
coordinate change on compactly supported tests; it is used before
placing the test in the radial Hilbert space.  The transformed test
$\Theta D_{\log b}G$ is supported strictly inside the unit disk.
The same scaling calculation makes $\Gamma$ independent of $b$.
Write $G_R=D_{-\log R}G$ for the outward rescaling of $G$.
For fixed support shapes and sufficiently large $R$,
Proposition~\ref{prop:radial-gap} gives
\begin{equation}
 \left|\Cov(F,G_R)-R^{-1}\Lambda(F)\Gamma(G)\right|
 \leq C R^{-3}\|F\|_2\|G\|_2.
 \label{eq:transfer-general-radial}
\end{equation}
Indeed the radius ratio after the preceding two normalizations is
$a_0/(bR)$; its first-order factor is exactly the one displayed.
The constants depend only on the fixed support enlargements, and not
on $R$.  All these assertions hold after square-group averaging.

\subsection{One lattice measurement for every bounded test}

Fix once and for all a bounded continuous test $W$ of finitely many
smeared fields in an annulus outside the unit disk, even and
square-invariant, with
\begin{equation}
 g:=\Gamma(W)\ne0.
 \label{eq:transfer-measuring-test}
\end{equation}
Such a test is obtained by reflecting the bounded test following
Proposition~\ref{prop:radial-gap}, whose support may be kept away
from zero.  Also fix a bounded continuous, even, square-invariant test
$V$ from the same smooth-bump construction, supported strictly inside
$[-1,1]^2$, with $\Lambda(V)\ne0$.  The spatial smears defining $V$
and $W$ are smooth, compactly supported, and fixed once and for all.
For an integer $t$, denote by $W_t$ and $V_t$ their lattice versions
at scale $t$.  Each smeared coordinate is formed by the normalized
sum $t^{-15/8}\sum_x\psi(x/t)\sigma_x$, including the fixed lattice
normalization from Proposition~\ref{prop:pure-spin}, using those same
fixed spatial functions at every scale and averaging over the square
group if necessary.  Their norms
are uniformly bounded and their joint laws, together with other fixed
smears, converge to their continuum counterparts.

Choose a fixed large integer $H$, to be specified after $L$, and define
for every $s$
\begin{equation}
 \ell_s(f)=\frac H g\Cov_0(f,W_{Hs}).
 \label{eq:transfer-measured-coefficient}
\end{equation}
This is a linear functional annihilating constants.  It depends on a
full-plane law and fixed bulk tests only.

We make the uniform meaning of this measurement explicit.  For any
sequence of functions $f_s$ with $\qnorm{f_s}\leq1$, choose their
midrange representatives, so $\|f_s\|_\infty\leq1$.
For any fixed finite list of large radius ratios, including $H$,
Lemma~\ref{lem:bounded-replacement} gives a common continuum proxy
$F$ along a subsequence, with $\|F\|_\infty\leq1$, for all the
required pairings.  Applying \eqref{eq:transfer-general-radial} at
ratio $H$ gives
\begin{equation}
 \ell_s(f_s)=\Lambda(F)+O(H^{-2})+o(1)+O(Ht).
 \label{eq:transfer-proxy-measurement}
\end{equation}
The $t$ is the chosen replacement error, the $o(1)$ is the subsequent
mesh error, and the constants are uniform in the sequence.  This is
an estimate for the pairings of a proxy, not a claim that $f_s$ itself
converges as a random field test.  The same proxy must be used for the
finite list, which is why the simultaneous assertion in
Lemma~\ref{lem:bounded-replacement} was retained.

Taking $t$ arbitrarily small after fixing $H$ proves a uniform bound
on $\ell_s$ for all sufficiently large $s$, independent of large $H$.
Likewise, with the fixed test $V_s$,
\[
 \ell_s(V_s)=\Lambda(V)+O(H^{-2})+o(1).
\]
For all sufficiently large fixed $H$ and then large $s$ this is
bounded away from zero by $|\Lambda(V)|/2$.  Therefore
\begin{equation}
 e_s=V_s/\ell_s(V_s)
 \label{eq:transfer-right-inverse}
\end{equation}
is a uniformly bounded right inverse.  Both constants are independent
of the subsequently prescribed error tolerance and of $L$.

\subsection{Boundary uniformity and compatibility across scales}

The measured coefficients and their bounded right inverses are now fixed.
It remains to establish the two operator estimates.  We first turn an
arbitrary boundary assignment into a bounded test across the annulus.
Let $\tau$ be any fixed
assignment on $\partial B_{Ls}$ and take the intermediate contour
$I_s=\partial B_{\lfloor Ls/2\rfloor}$.  Its marginal under that
boundary condition has a density $b_{s,\tau}$ relative to its
full-plane marginal satisfying
\begin{equation}
 K^{-1}\leq b_{s,\tau}\leq K,\qquad \E_0 b_{s,\tau}=1,
 \label{eq:transfer-boundary-density}
\end{equation}
by Proposition~\ref{prop:pure-buffered}.  The constant $K$ is
independent of $L,s,\tau$: the contour-to-boundary ratio is fixed.
The spatial Markov property gives the exact identity
\begin{equation}
 \E_0[f\mid\sigma_{\partial B_{Ls}}=\tau]-\E_0f
       =\Cov_0(f,b_{s,\tau}).
 \label{eq:transfer-boundary-pairing}
\end{equation}
When $f$ is even and square-invariant, we may average $b_{s,\tau}$
under those symmetries without changing this pairing or its bounds.

Consider any sequence of functions and boundary assignments in
\eqref{eq:transfer-boundary-pairing}.  Apply bounded replacement to
the inner support, to a fixed proportional band around $I_s$, and to
the measuring test at ratio $H$.  The enlargements around the inner
square have fixed size in units of $s$; the band enlargement has fixed
size in units of $Ls$ and stays a positive distance from zero in
those coordinates.  For $L$ large enough they are disjoint.  The
resulting proxies $F,G$ satisfy $\|F\|_\infty\leq1$ and
$\|G\|_\infty\leq K$.  Equation~\eqref{eq:transfer-general-radial}
therefore implies
\[
 L^2|\Cov(F,G_L)|
 \leq CL|\Lambda(F)|+C/L.
\]
Substituting \eqref{eq:transfer-proxy-measurement}, and letting the
replacement accuracy tend to zero after fixing $L,H$, proves
\begin{equation}
 \limsup_{s\to\infty}\sup_{\qnorm f\leq1}
 \left(\qnorm{A_sf}-CL|\ell_s(f)|\right)
 \leq C/L+CL/H^2.
 \label{eq:transfer-first-limsup}
\end{equation}
For clarity, the centered supremum in
\eqref{eq:transfer-boundary-pairing} bounds $\qnorm{A_sf}$ from
above.  Uniformity in \eqref{eq:transfer-first-limsup} follows by
contradiction: any failure supplies a sequence of bounded functions
and a maximizing boundary assignment, and the preceding common-proxy
argument applies to that sequence.  Thus no compactness of the
space of microscopic boundary assignments is being assumed.

The compatibility estimate uses an exact Markov identity before any
approximation.  Since $W_{HLs}$ is supported outside $B_{Ls}$,
\begin{align}
 \ell_{Ls}(A_sf)
 &=\frac{HL^2}{g}\Cov_0\bigl(
      \E_0[f\mid\sigma_{\partial B_{Ls}}],W_{HLs}\bigr)\notag\\
 &=\frac{HL^2}{g}\Cov_0(f,W_{HLs}).
 \label{eq:transfer-exact-tower}
\end{align}
The second equality is conditional independence of the inside and
outside given the boundary; it is not a formal composition of unrelated
conditional expectations.  Use a common proxy for $f$ in its pairings
at ratios $H$ and $HL$.  Equation~\eqref{eq:transfer-general-radial}
then gives respectively
\begin{align*}
 \ell_s(f)&=\Lambda(F)+O(H^{-2})+o(1)+O(Ht),\\
 \ell_{Ls}(A_sf)&=L\Lambda(F)
                  +O((LH^2)^{-1})+o(1)+O(HL^2t).
\end{align*}
After fixing $L,H$ and taking the accuracies in their stated order,
we obtain
\begin{equation}
 \limsup_{s\to\infty}\sup_{\qnorm f\leq1}
       |\ell_{Ls}(A_sf)-L\ell_s(f)|\leq CL/H^2.
 \label{eq:transfer-second-limsup}
\end{equation}
The constants in these two displays depend on the fixed chart shapes
and measuring test, and are independent of $L,H$ once both are large.

\begin{proof}[Proof of Theorem~\ref{thm:transfer}]
Choose $L$ large enough that the term $C/L$ in
\eqref{eq:transfer-first-limsup} is less than $a/4$ and the fixed
support enlargements are disjoint.  Next choose an integer $H>L$
large enough that all terms $CL/H^2$ are less than $a/4$, and that
the lower bound needed in \eqref{eq:transfer-right-inverse} holds.
Choose every replacement accuracy small compared with the now fixed
amplification factors $H,L^2,HL^2$.  In each replacement construction,
choose the window, its finite subdivision, the observation strength,
and the coarse partition before taking the mesh small.  Finally choose
$s_0$ so that the uniform upper limits
\eqref{eq:transfer-first-limsup} and
\eqref{eq:transfer-second-limsup} are within their remaining error
margins, and so that the norm and right-inverse bounds hold for every
$s\geq s_0$.  Homogeneity gives the estimates for arbitrary $f$.

All square contours and cell boundaries used here are specified by
integer rounding.  They remain in the same fixed buffered charts;
their rescaled locations converge to the prescribed contours and
cells.  Changes of cell indicators are confined to shrinking
neighborhoods of their finitely many polygonal edges; their normalized
lattice counting measure tends to zero.  The uniform $L^2$ bounds on
fixed-order rescaled correlation kernels and Cauchy--Schwarz therefore
make the resulting smearing errors tend to zero.  Spatial test functions
and their finite menus are fixed before the mesh limit.  Only these
specified geometric roundings and the convergent cell-normalization
scalars depend on the mesh, not arbitrary varying Riemann-integrable
spatial functions.  The separate line-continuity estimate controls
the auxiliary square contours.  Consequently changing an auxiliary contour by a
bounded number of lattice layers affects none of the limiting
estimates or their uniform thresholds.  The conditioning convention in
\eqref{eq:transfer-operator}, however, is fixed throughout, so the
Markov identity \eqref{eq:transfer-exact-tower} is exact at every scale.
\end{proof}

For the analytic construction in the next section one may complexify
these spaces.  The complex quotient norm is
$\qnorm f=\inf_{c\in\C}\|f-c\|_\infty$; its linear mean-zero
representative has norm at most $2\qnorm f$.  Extend $A_s$ and
$\ell_s$ complex linearly.  Applying the real estimates to real and
imaginary parts changes universal constants by at most a factor two,
and the arbitrarily small error $a$ can be chosen smaller beforehand.
Thus the same conclusions, with harmless fixed constants, hold for
the complex spaces used to justify analyticity.
